% Zone „Kennst du schon" – Kreis, Kl. 8 – Nr. 1–9
\setcounter{aufgabe}{0}
\zonekopf

\begin{aufgabe}{Ich kann mit Kommazahlen malnehmen und teilen und das Ergebnis runden.}
Rechne im Kopf.
\begin{teilezwei}
\tz $0{,}5 \cdot 8 =$ \leerfeld & \tz $2{,}5 \cdot 4 =$ \leerfeld \\
\anweisung{Rechne mit dem Taschenrechner. Runde auf eine Stelle nach dem Komma.}
\tz $3{,}2 \cdot 1{,}5 \approx$ \leerfeld & \tz $4{,}37 \cdot 2{,}6 \approx$ \leerfeld \\
\tz $25 : 3{,}2 \approx$ \leerfeld & \tz $3{,}96 \cdot 5{,}05 \approx$ \leerfeld \\
\anweisung{Rechne ohne Zwischenrunden. Runde erst das Endergebnis auf eine Stelle.}
\tz $6{,}35 \cdot 1{,}5 \cdot 4 \approx$ \leerfeld & \\
\end{teilezwei}
\end{aufgabe}

\begin{aufgabe}{Ich kann Umfang und Flächeninhalt unterscheiden und mit der passenden Einheit angeben.}
Berechne den Umfang $u$ und den Flächeninhalt $A$. Schreibe die Einheit dazu.
\begin{geruest}
\gz{a}{Rechteck mit den Seiten $3\,\text{cm}$ und $5\,\text{cm}$}{\feld{u}\quad\feld{A}}
\gz{b}{Quadrat mit der Seite $4\,\text{m}$}{\feld{u}\quad\feld{A}}
\gz{c}{Rechteck mit den Seiten $2{,}5\,\text{m}$ und $4\,\text{m}$}{\feld{u}\quad\feld{A}}
\anweisung{Gib Umfang und Flächeninhalt in m und m² an.}
\gz{d}{Rechteck mit den Seiten $50\,\text{cm}$ und $2\,\text{m}$}{\feld{u}\quad\feld{A}}
\end{geruest}
\end{aufgabe}

\begin{aufgabe}{Ich kann eine Formel nach einer anderen Größe umstellen.}
Stelle die Formel nach der Größe in Klammern um.
\begin{gleichungsraster}[1]
\gl{u = 4 \cdot a \quad (a)} & \gl{A = a \cdot b \quad (b)} \\
\gl{V = a \cdot b \cdot c \quad (c)} & \\
\anweisung{Setze ein und berechne die gesuchte Größe.}
\gl[2]{u = 4 \cdot a;\ u = 18{,}4\,\text{cm}} & \gl[2]{V = a \cdot b \cdot c;\ V = 60\,\text{cm}^3,\ a = 5\,\text{cm},\ b = 4\,\text{cm}} \\
\end{gleichungsraster}
\end{aufgabe}

\begin{aufgabe}{Ich kann Zahlen quadrieren und die Wurzel ziehen.}
Rechne im Kopf.
\begin{teilezwei}
\tz $6^2 =$ \leerfeld & \tz $\sqrt{49} =$ \leerfeld \\
\anweisung{Rechne mit dem Taschenrechner. Runde auf eine Stelle nach dem Komma, wo es nötig ist.}
\tz $1{,}5^2 =$ \leerfeld & \tz $0{,}5^2 =$ \leerfeld \\
\tz $\sqrt{12{,}25} =$ \leerfeld & \tz $\sqrt{30} \approx$ \leerfeld \\
\anweisung{Welche positive Zahl ergibt quadriert die angegebene Zahl?}
\tz $\square^2 = 36$: \leerfeld & \tz $\square^2 = 2{,}56$: \leerfeld \\
\end{teilezwei}
\end{aufgabe}

\begin{aufgabe}{Ich finde den Fehler beim Quadrieren und Wurzelziehen.}
Leo hat so gerechnet. Finde den Fehler, benenne ihn und korrigiere die Rechnung.
\begin{teile}
\teil \rechnung{3{,}5^2 &= 2 \cdot 3{,}5 \\ &= 7}
\teil \rechnung{\sqrt{64} &= 64 : 2 \\ &= 32}
\end{teile}
\end{aufgabe}

\begin{aufgabe}{Ich kann Quadrat und Doppeltes auseinanderhalten.}
Berechne beides.
\begin{teilezwei}
\tz $4^2 =$ \leerfeld & \tz $2 \cdot 4 =$ \leerfeld \\
\tz $2{,}5^2 =$ \leerfeld & \tz $2 \cdot 2{,}5 =$ \leerfeld \\
\tz $\sqrt{36} =$ \leerfeld & \tz $36 : 2 =$ \leerfeld \\
\tz $0{,}3^2 =$ \leerfeld & \tz $2 \cdot 0{,}3 =$ \leerfeld \\
\end{teilezwei}
\end{aufgabe}

\begin{aufgabe}{Ich kann Winkel messen und zum Vollwinkel ergänzen.}
Gib den Winkel in Grad an.
\begin{teilezwei}
\tz Vollwinkel: \leerfeld[°] & \tz rechter Winkel: \leerfeld[°] \\
\tz gestreckter Winkel: \leerfeld[°] & \tz gemessener Winkel unten: \leerfeld[°] \\
\end{teilezwei}

\winkel[4]{65}{}

\anweisung{Wie viel Grad fehlen bis zum Vollwinkel?}
\begin{teilezwei}
\tz $130^\circ$ + \leerfeld[°] $= 360^\circ$ & \tz $275^\circ$ + \leerfeld[°] $= 360^\circ$ \\
\end{teilezwei}
\end{aufgabe}

\begin{aufgabe}{Ich kann einen Winkel mit dem Geodreieck zeichnen.}
Trage am Scheitel $S$ den Winkel an.
\begin{teile}
\teil $40^\circ$
\end{teile}
\winkelstrahl[6]
\begin{teile}
\teil $75^\circ$
\end{teile}
\winkelstrahl[6]
\begin{teile}
\teil $100^\circ$
\end{teile}
\winkelstrahl[6]
\end{aufgabe}

\begin{aufgabe}{Ich kann einen Anteil als Bruch und in Prozent angeben.}
Gib den Anteil als gekürzten Bruch und in Prozent an.
\begin{teile}
\teil 5 von 20 \hfill \leerfeld \quad \leerfeld[\%]
\teil 30 von 60 \hfill \leerfeld \quad \leerfeld[\%]
\teil 42 von 120 \hfill \leerfeld \quad \leerfeld[\%]
\teil 50 von 360, Prozent auf eine Stelle gerundet \hfill \leerfeld \quad \leerfeld[\%]
\teil Von 80 Kindern fahren 20 mit dem Rad. Welcher Anteil fährt nicht mit dem Rad? \\ \leerfeld \quad \leerfeld[\%]
\end{teile}
\end{aufgabe}
